\textbf{2-1}\quad 证明空间反演算符 $\hat{\Pi}[\hat{\Pi}\psi(x)=\psi(-x)]$ 是厄米算符. 指出在什么条件下, $\hat{p}=-i\hbar\frac{d}{dx}$ 是厄米算符.

\vfill

\newpage
\vfill

\noindent
\textbf{2-2}\quad 动量在径向方向的分量定义为 $\hat{p}_r = \frac{1}{2}\left(\hat{\pmb{p}}\cdot \frac{\pmb{r}}{r} +\frac{\pmb{r}}{r}\cdot \hat{\pmb{p}}\right)$ .求出 $\hat{p}_r$ 在球坐标中的表达式.

\vfill

\newpage
\vfill

\noindent
\textbf{2-3}\quad 证明 $[\hat{p}_x, f(x)] = -\mathrm{i}\hbar \frac{\partial}{\partial x} f(x), [x, f(\hat{p}_x)] = \mathrm{i}\hbar \frac{\partial}{\partial \hat{p}_x} f(\hat{p}_x)$

\vfill

\newpage
\vfill

\noindent
\textbf{2-4}\quad 设算符 $\hat{A}$ 满足条件 $\hat{A}^2 = 1$ ，证明 $\mathrm{e}^{\mathrm{i}\alpha \hat{A}} = \cos \alpha + \mathrm{i}\sin \alpha \hat{A}$ ，其中 $\alpha$ 为实常数。

\vfill

\newpage
\vfill

\noindent
\textbf{2-5}\quad 设算符 $\hat{K} = \hat{L}\hat{M}$ , $\hat{L}\hat{M} - \hat{M}\hat{L} = 1$ , 又设 $\varphi$ 为 $\hat{K}$ 的本征矢, 相应本征值为 $\lambda$ . 证明 $u \equiv \hat{L}\varphi$ 和 $v \equiv \hat{M}\varphi$ 也是 $\hat{K}$ 的本征矢, 并求出相应的本征值.

\vfill

\newpage
\vfill

\noindent
\textbf{2-6}\quad 粒子作一维运动， $\hat{H} = \frac{\hat{p}^2}{2\mu} + V(x)$ ，定态波函数为 $|n\rangle: \hat{H}|n\rangle = E_n |n\rangle, n = 1, 2, \cdots$ 

(1) 证明
\begin{equation}\tag{1}
\langle n \left| \hat {p} \right| m \rangle = a _ {n m} \langle n | x | m \rangle
\end{equation}

并求系数 $a_{nm}$ . 

(2) 利用式(1)推导求和公式

\begin{equation}\tag{2}
\sum_ {n} (E _ {n} - E _ {m}) ^ {2} \left| \left\langle n \right| x \mid m \right\rangle \Big | ^ {2} = \frac {\hbar^ {2}}{\mu^ {2}} \left\langle m \right| p ^ {2} \left| m \right\rangle
\end{equation}

(3) 证明

\begin{equation}\tag{3}
\sum_ {n} \left(E _ {n} - E _ {m}\right) \left| \left\langle n \mid x \mid m \right\rangle \right| ^ {2} = \frac {\hbar^ {2}}{2 \mu}
\end{equation}

\vfill

\newpage
\vfill

\noindent
\textbf{2-7}\quad 设 $\hat{F}$ 为厄米算符, 证明在能量表象中下式成立:

$$
\sum_ {n} (E _ {n} - E _ {k}) \left| F _ {n k} \right| ^ {2} = \frac {1}{2} \left\langle k \left| [ \hat {F}, [ \hat {H}, \hat {F} ] ] \right| k \right\rangle
$$

\vfill

\newpage
\vfill

\noindent
\textbf{2-8}\quad 有一量子力学体系, 哈密顿量 $\hat{H}$ 的本征值与本征矢量分别为 $E_{n}$ 与 $\left|n\right\rangle: \hat{H}\left|n\right\rangle = E_{n}\left|n\right\rangle$ . 设 $\hat{F}$ 为任一算符 $\hat{F} = \hat{F}(x, \hat{p})$ , 试证明

\vfill

\newpage
\vfill

\noindent
\textbf{2-9}\quad 已知 $Y_{lm}(\theta, \varphi)$ 是 $\hat{L}^2$ 和 $\hat{L}_z$ 的共同本征函数, 本征值分别为 $l(l + 1)\hbar^2$ 和 $m\hbar$ . 令 $\hat{L}_{\pm} = \hat{L}_x \pm \mathrm{i}\hat{L}_y$ . 

(1) 证明 $\hat{L}_{\pm}Y_{lm}(\theta, \varphi)$ 仍是 $\hat{L}^2$ 和 $\hat{L}_z$ 的共同本征函数, 求出它们的本征值；

(2) 推导公式 $\hat{L}_{\pm}Y_{lm}(\theta, \varphi) = \sqrt{l(l + 1) - m(m \pm 1)}\hbar Y_{lm \pm 1}(\theta, \varphi)$ .

\vfill

\newpage
\vfill

\noindent
\textbf{2-10}\quad 证明 $\mathbf{e}^{\hat{A}}\hat{B}\mathbf{e}^{-\hat{A}} = \hat{B} + [\hat{A},\hat{B}] + \frac{1}{2!} [\hat{A},[\hat{A},\hat{B}]] + \frac{1}{3!} [\hat{A},[\hat{A},[\hat{A},\hat{B}]]] + \dots .$

\vfill

\newpage
\vfill

\noindent
\textbf{2-11}\quad 设算符 $\hat{A}$ 与 $\hat{B}$ 同它们的对易关系式 $[\hat{A},\hat{B}]$ 都对易, 证明

$$
[ \hat {A}, \hat {B} ^ {n} ] = n \hat {B} ^ {n - 1} [ \hat {A}, \hat {B} ]
$$

$$
\mathbf {e} ^ {\hat {A} + \hat {B}} = \mathbf {e} ^ {\hat {A}} \mathbf {e} ^ {\hat {B}} \mathbf {e} ^ {- \frac {1}{2} [ \hat {A}, \hat {B} ]} \quad \text {或} \quad \mathbf {e} ^ {\hat {A}} \mathbf {e} ^ {\hat {B}} = \mathbf {e} ^ {\hat {A} + \hat {B} + \frac {1}{2} [ \hat {A}, \hat {B} ]}
$$

\vfill

\newpage
\vfill

\noindent
\textbf{2-12}\quad 设 $\hat{L}$ 为轨道角动量算符. 已知 $\hat{L}^2$ 与 $\hat{L}_z$ 的共同本征函数为 $Y_{lm}(\theta, \varphi)$ . 证明 $\mathrm{e}^{\frac{1}{\mathrm{i}\hbar}\hat{L}_y\frac{\pi}{2}} Y_{lm}(\theta, \varphi)$ 为 $\hat{L}^2$ 和 $\hat{L}_x$ 的共同本征函数, 并求出相应的本征值.

\vfill

\newpage
\vfill

\noindent
\textbf{2-13}\quad 设 $\hat{p}$ 为 $x$ 方向的动量算符，满足对易关系 $[x, \hat{p}] = \mathrm{i}\hbar$ 。求 

(1) $\mathrm{e}^{\mathrm{i}a\hat{p}} x \mathrm{e}^{-\mathrm{i}a\hat{p}} = ?$ 

(2) $\mathrm{e}^{-\mathrm{i}b\hat{p} x} [x, \mathrm{e}^{\mathrm{i}b\hat{p} x}] = ?$ 其中 $a, b$ 为常数。

\vfill

\newpage
\vfill

\noindent
\textbf{2-14}\quad 设粒子处于状态 $Y_{lm}(\theta, \varphi)$ , 求轨道角动量 $x$ 分量及 $y$ 分量平均值 $\overline{L}_x$ 与 $\overline{L}_y$ , 以及 $\overline{(\Delta L_x)^2}$ 与 $\overline{(\Delta L_y)^2}$ .

\vfill

\newpage
\vfill

\noindent
\textbf{2-15}\quad 已知角动量算符的三个分量 $\hat{J}_x, \hat{J}_y, \hat{J}_z$ 满足对易关系

$$
[ \hat {J} _ {x}, \hat {J} _ {y} ] = \mathrm{i} \hbar \hat {J} _ {z}, [ \hat {J} _ {y}, \hat {J} _ {z} ] = \mathrm{i} \hbar \hat {J} _ {x}, [ \hat {J} _ {z}, \hat {J} _ {x} ] = \mathrm{i} \hbar \hat {J} _ {y}
$$

定义: $\hat{J}^2 = \hat{J}_x^2 + \hat{J}_y^2 + \hat{J}_z^2$ , $\hat{J}_{\pm} = \hat{J}_x \pm i\hat{J}_y$ . 

(1) 求对易关系 $[\hat{J}^2, \hat{J}_{\pm}]$ , $[\hat{J}_z, \hat{J}_{\pm}]$ , $[\hat{J}_{+}, \hat{J}_{-}]$ ;

(2)若 $\hat{J}^2$ 和 $\hat{J}_z$ 的共同本征函数为 $\psi_{jm}$ , 其中 $j$ 和 $m$ 为相应的量子数, 证明 $\hat{J}_{\pm} \psi_{jm}$ 也是 $\hat{J}^2$ 和 $\hat{J}_z$ 的共同本征函数, 并求出相应的本征值.

\vfill

\newpage
\vfill

\noindent
\textbf{2-16}\quad 设 $\hat{J}_x, \hat{J}_y, \hat{J}_z$ 为角动量算符， $\hat{J}_{\pm} = \hat{J}_x \pm \mathrm{i}\hat{J}_y$ 。算符 $\hat{V}_{+}$ 与 $\hat{J}_z, \hat{J}_{+}$ 满足对易关系： $[\hat{J}_{+}, \hat{V}_{+}] = 0, [\hat{J}_z, \hat{V}_{+}] = \hbar \hat{V}_{+}$ 。证明 $\hat{V}_{+}|jj\rangle = c|j + 1, j + 1\rangle$ ，其中 $c$ 为常数， $|jm\rangle$ 为 $\hat{J}^2$ 与 $\hat{J}_z$ 的共同本征函数。

\vfill

\newpage
\vfill

\noindent
\textbf{2-17}\quad 令 $\hat{p}_{+} = \hat{p}_{x} + \mathrm{i}\hat{p}_{y},\hat{L}_{+} = \hat{L}_{x} + \mathrm{i}\hat{L}_{y}$ . 

(1) 计算 $[\hat{L}_x,\hat{p}_+],[\hat{L}_y,\hat{p}_+],[\hat{L}_z,\hat{p}_+]$ . 证明

$$
[ \hat {L} ^ {2}, \hat {p} _ {+} ] = 2 \hbar (\hat {p} _ {+} \hat {L} _ {z} - \hat {p} _ {z} \hat {L} _ {+}) + 2 \hbar^ {2} \hat {p} _ {+}
$$

(2) 已知 $\hat{L}_z\Phi_m = m\hbar \Phi_m$ . 证明 $\hat{p}_+ \Phi_m$ 仍是 $\hat{L}_z$ 的本征态, 并求出它的本征值. (3) 对于自由粒子体系, 已知其哈密顿算符 $\hat{H}$ 同 $\tilde{L}^2$ 与 $\hat{L}_z$ 有共同的本征函数完备系 $\psi_{klm} = R_{kl}(r)Y_{lm}(\theta, \varphi)$ , 相应的本征值分别是 $E_k = \hbar^2 k^2 / 2\mu, l(l + 1)\hbar, m\hbar$ . 试证明 $\hat{p}_+ \psi_{kll}$ 仍是 $\hat{H}, \hat{L}^2$ 与 $\hat{L}_z$ 的共同本征函数, 并求出相应的本征值.

\vfill

\newpage
\vfill

\noindent
\textbf{2-18}\quad 设算符 $\hat{H}$ 具有连续本征值, 其本征函数 $u(x, \omega)$ 构成正交完备系. 求方程
\begin{equation}
(\hat{H} - \omega_{0}) V(x) = F(x)\tag{1}
\end{equation}
的解, 其中 $F(x)$ 为已知函数, $\omega_0$ 为某个特定的本征值.

\vfill

\newpage
\vfill

\noindent
\textbf{2-19}\quad 定义平移算符 $U_{x}(a)$ , 它对波函数 $\psi(x)$ 的作用是
$$
U _ {x} (a) \psi (x) = \psi (x - a)
$$
其中 $a$ 为实数. 

(1) 证明 $U_x(a) = \mathrm{e}^{-\mathrm{i}a\hat{p}_x / \hbar}$ ; 

(2) 证明 $U_x(a)$ 为幺正算符.

\vfill

\newpage
\vfill

\noindent
\textbf{2-20}\quad 一维谐振子处于定态 $\psi_{n}$ , 计算 $\Delta x\Delta p$ , 检验测不准关系

\vfill

\newpage
\vfill

\noindent
\textbf{2-21}\quad t=0 时一维自由运动粒子的归一化波函数为
\begin{equation}
\psi (x, 0) = \left(2 \pi a ^ {2}\right) ^ {- 1 / 4} \exp \left[ i k _ {0} \left(x - x _ {0}\right) - \frac {\left(x - x _ {0}\right) ^ {2}}{(2 a) ^ {2}} \right]
\end{equation}
其中 $a, k_0$ 与 $x_0$ 均为正实数。

(1) 求 $t = 0$ 时粒子的坐标概率分布函数与坐标分布宽度 $\Delta x = \sqrt{x^2 - (\bar{x})^2}$ ; 

(2) 求 $t = 0$ 时粒子的动量概率分布函数与动量分布宽度 $\Delta p = \sqrt{p^2 - (\bar{p})^2}$ , 并检验坐标与动量的测不准关系; 

(3) 求 $t > 0$ 时粒子的波函数 $\psi(x, t)$ , 坐标 $x$ 的平均值 $\overline{x(t)}$ 及坐标的概率分布函数; 

(4) 求 $t > 0$ 时粒子动量的平均值 $\overline{p(t)}$ 及动量的概率分布函数.

\vfill

\newpage
\vfill

\noindent
\textbf{2-22}\quad 已知束缚态波函数为 $\psi(x)$ , 求动量 $p$ 与动能 $T = p^2 / 2\mu$ 的概率分布函数的表达式. 对一维谐振子基态, 波函数为 $\psi(x) = \sqrt{\frac{\alpha}{\sqrt{\pi}}} e^{-\alpha^2 x^2 / 2}, \alpha = \sqrt{\mu \omega / \hbar}$ . 算出动量 $p$ 与动能 $T$ 的概率分布函数, 并算出动能平均值.

\vfill

\newpage
\vfill

\noindent
\textbf{2-23}\quad 一维谐振子能量的本征值与本征函数为
\begin{equation}\tag{1}
E _ {n} = \left(n + \frac {1}{2}\right) \hbar \omega , \psi_ {n} (x) = N _ {n} \mathrm{e} ^ {- \alpha^ {2} x ^ {2} / 2} H _ {n} (\alpha x)
\end{equation}

\begin{equation}\tag{2}
N _ {n} = \sqrt {\frac {\alpha}{\sqrt {\pi} 2 ^ {n} n !}}, \alpha = \sqrt {\frac {\mu \omega}{\hbar}}, n = 0, 1, 2, \dots
\end{equation}

(1) 由厄米多项式 $H_{n}(z)$ 的递推关系
\begin{equation}\tag{3}
z H _ {n} (z) = \frac {1}{2} H _ {n + 1} (z) + n H _ {n - 1} (z)
\end{equation}
\begin{equation}\tag{4}
H _ {n} ^ {\prime} (z) = 2 n H _ {n - 1} (z)
\end{equation}
导出 $x \psi_{n}(x)$ 和 $\mathrm{d} \psi_{n}(x) / \mathrm{d} x$ 满足的递推关系。

(2) 求出 $\psi_{n}(x)$ 态上坐标和动量的平均值 $\overline{x}$ 和 $\overline{p}$ . 

(3) 证明谐振子零点能 $E_{0} = \frac{\hbar \omega}{2}$ 是测不准关系 $\Delta x \Delta p \geqslant \frac{\hbar}{2}$ 的直接结果. 

(4) 求出 $\psi_{n}(x)$ 态上动能与势能的平均值 $\overline{T}$ 和 $\overline{V}$ , 并找出它们之间的关系.

\vfill

\newpage
\vfill

\noindent
\textbf{2-24}\quad 质量为 $\mu$ 的粒子在外场作用下作一维运动 $(- \infty < x < + \infty)$ . 已知当其处于束缚态 $\psi_{1}(x)$ 时, 动能平均值等于 $E_{1}$ , 并已知 $\psi_{1}(x)$ 是实函数. 试求当粒子处于态 $\psi_{2}(x) = \psi_{1}(x)e^{ikx}$ ( $k$ 为实数)时动量平均值 $\overline{p}$ 与动能平均值 $\overline{T}$ .

\vfill

\newpage
\vfill

\noindent
\textbf{2-25}\quad 一维谐振子哈密顿算符(取 $\hbar = \mu = \omega = 1$ ) 为
\begin{equation}\tag{1}
\hat {H} = \frac {1}{2} (x ^ {2} + \hat {p} ^ {2})
\end{equation}
其本征值与本征函数为 $E_{n}$ 与 $\psi_{n}(x)$ . 已知 $\psi_{n}(x)$ 为实函数, 宇称为 $(-1)^{n}$ . 请写出 $E_{n}$ 的具体形式. 已知 $t = 0$ 时谐振子的波函数为
\begin{equation}\tag{2}
\psi (0) = \frac {1}{2} \left(\sqrt {2} \psi_ {0} + \psi_ {1} + \psi_ {2}\right)
\end{equation}
求任意 $t$ 时刻 $x$ 与 $\hat{p}$ 的平均值 $\overline{x(t)}$ 与 $\overline{p(t)}$ .

\vfill

\newpage
\vfill

\noindent
\textbf{2-26}\quad 证明在宽度为 $a$ 的一维无限深方势阱中的定态能量 $E > \hbar^2 / 2\mu a^2$

\vfill

\newpage
\vfill

\noindent
\textbf{2-27}\quad (1)如厄米算符 $\hat{A}$ 对任何态矢量 $|u\rangle$ , 有 $\langle u|\hat{A}|u\rangle > 0$ , 则称 $\hat{A}$ 是正定算符. 求证算符 $\hat{A} = |a\rangle\langle a|$ 是厄米正定算符. 

(2) 如 $\hat{A}$ 是任一线性算符, 证明 $\hat{A}^{+}\hat{A}$ 是厄米正定算符, 它的迹等于 $\hat{A}$ 在任意表象中的矩阵元的模平方之和, 试推导, 当且仅当 $\hat{A} = 0$ 时, $\operatorname{tr}(\hat{A}^{+}\hat{A}) = 0$ 才成立.

\vfill

\newpage
\vfill

\noindent
\textbf{2-28}\quad 设归一化波函数 $|\psi\rangle$ 满足薛定谔方程 $\mathrm{i}\hbar \frac{\partial}{\partial t} |\psi \rangle = \hat{H} |\psi \rangle$ 。定义密度算符(矩阵)为 $\rho = |\psi\rangle \langle \psi|$ 。

(1)证明任意力学量 $\hat{F}$ 在 $|\psi\rangle$ 态中的平均值可表示为 $\operatorname{tr}(\hat{F}\rho)$ ; 

(2)求出 $\rho$ 的本征值；

(3)导出 $\rho$ 随时间变化的方程。

\vfill

\newpage
\vfill

\noindent
\textbf{2-29}\quad 已知一量子体系, 除了能量之外还包括另外三个可观察量 $\hat{P}, \hat{Q}$ 与 $\hat{R}$ . 设该体系只有两个能量本征态 $|1\rangle$ 与 $|2\rangle$ , 它们不一定是 $\hat{P}, \hat{Q}$ 与 $\hat{R}$ 的本征态. 基于以下各组"实验数据", 尽可能多地定出 $\hat{P}, \hat{Q}$ 与 $\hat{R}$ 的本征值(有一组数据是非物理的): 

(1) $\langle 1|\hat{P}|1\rangle = 1/2$ , $\langle 1|\hat{P}^2|1\rangle = 1/4$ ; 

(2) $\langle 1|\hat{Q}|1\rangle = 1/2$ , $\langle 1|\hat{Q}^2|1\rangle = 1/6$ ; 

(3) $\langle 1|\hat{R}|1\rangle = 1$ , $\langle 1|\hat{R}^2|1\rangle = 5/4$ , $\langle 1|\hat{R}^3|1\rangle = 7/4$ .

\vfill

\newpage
\vfill

\noindent
\textbf{2-30}\quad 已知一体系哈密顿量 $\hat{H}$ 不含 $t$ , $\hat{H}$ 具有非简并本征值 $E_{\lambda}, \lambda = 1, 2, \cdots$ , 相应本征态矢为 $|\psi_{\lambda}\rangle: \hat{H}|\psi_{\lambda}\rangle = E_{\lambda}|\psi_{\lambda}\rangle$ . 又已知另一可观察量 $\hat{A}$ 的非简并本征值与本征态矢为 $a_{n}$ 与 $|\phi_{n}\rangle: \hat{A}|\phi_{n}\rangle = a_{n}|\phi_{n}\rangle, n = 1, 2, \cdots$ . 设体系初态为 $|\psi_{\lambda}\rangle$ , 问

(1)在初态 $|\psi_{\lambda}\rangle$ 下对 $\hat{A}$ 测量时, 测得 $\hat{A}$ 的平均值是什么? 这一测量给出 $\hat{A}$ 的值为 $a_{m}$ 的概率有多大? 

(2)如果上次对 $\hat{A}$ 测量得 $a_{m}$ , 经过时间间隔 $t$ 以后再次测量 $\hat{A}$ 仍得 $a_{m}$ 的概率有多大?

\vfill

\newpage
\vfill

\noindent
\textbf{2-31}\quad 计算 $\hat{\pmb{p}}\times \hat{\pmb{L}} +\hat{\pmb{L}}\times \hat{\pmb{p}} = ?$

\vfill

\newpage
\vfill

\noindent
\textbf{2-32}\quad 定义角动量升降算符 $\hat{J}_{\pm} = \hat{J}_x \pm i\hat{J}_y$ ，

(1)证明算符 $\hat{J}_{+}\hat{J}_{-}$ 与 $\hat{J}_{-}\hat{J}_{+}$ 的厄米性，并求出它们的本征态与本征值；

(2)若力学量算符 $\hat{F}$ 满足对易关系 $[\hat{F}, \hat{J}_{u}] = 0, \mu = x, y, z$ ，试证明 $\hat{F}$ 在 $\hat{J}^2, \hat{J}_z$ 共同本征态上的平均值与磁量子数无关。

\vfill

\newpage
\vfill

\noindent
\textbf{2-33}\quad 证明力学量 $\hat{A}$ (不显含 $t$ ) 的平均值对时间的二次微商为
$$
\frac {\mathrm{d} ^ {2} \overline {{A}}}{\mathrm{d} t ^ {2}} = - \frac {1}{\hbar^ {2}} \overline {{\left[ \left[ \hat {A} , \hat {H} \right] , \hat {H} \right]}}
$$
其中 $\hat{H}$ 为该体系的哈密顿量.

\vfill

\newpage
\vfill

\noindent
\textbf{2-34}\quad 证明在一维束缚态问题中,位能在基态的平均值满足如下关系:
$$
\left\langle V \right\rangle_ {0} = E _ {0} - \frac {m}{2 \hbar^ {2}} \sum_ {n} (E _ {n} - E _ {0}) ^ {2} \left| \left\langle n \right| x \left| 0 \right\rangle \right| ^ {2}
$$
再证明 $\left\langle V\right\rangle_0\leqslant (5E_0 - E_1) / 4,\quad E_0\leqslant E_1\leqslant E_2\leqslant \dots .$

\vfill

\newpage
\vfill

\noindent
\textbf{2-35}\quad 能量为 $E$ 的一束粒子穿过小孔, 投射到距离小孔距离为 $L$ 的屏上. 用不确定原理, 证明通过不断减小孔径来达到减小屏上斑点直径是不可行的. 估算使屏上斑点直径最小时小孔的直径.

\vfill

\newpage
\vfill

\noindent
\textbf{2-36}\quad 设体系的能量本征方程为 $\hat{H}\left|n\right\rangle = E_n\left|n\right\rangle, \left\langle m\mid n\right\rangle = \delta_{mn}, E_0 \leqslant E_1 \leqslant E_2 \cdots$ . 

(1) 取 $\left|\psi_0\right\rangle$ 为归一化基态试探态矢. 令 $E = \left\langle \psi_0\mid \hat{H}\mid \psi_0\right\rangle, \varepsilon = 1 - \left|\left\langle 0\mid \psi_0\right\rangle\right|^2$ , 证明 $E - E_0 > (E_1 - E_0)\varepsilon$ . 

(2) 若只知 $\hat{H}$ 最低的两个本征态矢 $\left|0\right\rangle$ 与 $\left|1\right\rangle$ , 试从任意归一化态矢出发, 构造第二激发态的试探态矢, 并求出该激发态能量上限.

\vfill

\newpage
\vfill

\noindent
\textbf{2-37}\quad 已知轨道角动量 $\hat{L}$ 在 n 方向上的分量为
$$
\hat {L} _ {n} = \hat {\boldsymbol {L}} \cdot \boldsymbol {n} = \hat {L} _ {x} \sin \alpha \cos \beta + \hat {L} _ {y} \sin \alpha \sin \beta + \hat {L} _ {z} \cos \alpha
$$
其中 $\alpha, \beta$ 为已知的方位角. 求在算符 $\hat{L}^2$ 与 $\hat{L}_z$ 的共同本征态 $|lm\rangle$ 上算符 $\hat{L}_n$ 和 $\hat{L}_n^2$ 的平均值.

\vfill

\newpage
\vfill

\noindent
\textbf{2-38}\quad 证明在 $\hat{H} = \frac{\hat{p}^2}{2\mu} + V(x)$ 的束缚定态 $\psi_n(x)$ 上, 动量 $\hat{p}$ 与作用力 $\hat{F}$ 的平均值为 0.

\vfill

\newpage
\vfill

\noindent
\textbf{2-39}\quad 设体系的哈密顿量 $\hat{H}$ 同力学量 $\hat{A}$ 满足反对易关系 $\hat{H}\hat{A} +\hat{A}\hat{H} = 0$ .设 $\psi$ 是 $\hat{H}$ 的本征值为 $E(\neq 0)$ 的本征态，

(1)证明 $\hat{A}\psi$ 是 $\hat{H}$ 的本征值为- $E$ 的本征态；

(2)求 $\hat{A}$ 在 $\psi$ 态上的平均值.

\vfill

\newpage
\vfill

\noindent
\textbf{2-40}\quad 已知力学量 $\hat{A}$ 与 $\hat{B}$ 的本征值分别为 $a_{n}$ 与 $b_{n}$ . 在 $\left|\psi\right\rangle$ 态上先测量 $\hat{A}$ 得 $a_{n}$ , 后测量 $\hat{B}$ 得 $b_{n}$ 的概率为 $P(a_{n}, b_{n})$ ; 先测量 $\hat{B}$ 得 $b_{n}$ , 后测量 $\hat{A}$ 得 $a_{n}$ 的概率为 $P(b_{n}, a_{n})$ . 问 $P(a_{n}, b_{n}) = P(b_{n}, a_{n})$ 的条件是什么?

\vfill

\newpage
\vfill

\noindent
\textbf{2-41}\quad 已知可观察量 $A$ 的算符 $\hat{A}$ 有两个本征函数 $\phi_1, \phi_2$ ，本征值分别为 $a_1, a_2$ ；观察量 $B$ 的算符 $\hat{B}$ 有两个本征函数 $\chi_1, \chi_2$ ，本征值分别为 $b_1, b_2$ 。两种本征态有如下关系：
$$
\phi_ {1} = \frac {2 \chi_ {1} + 3 \chi_ {2}}{\sqrt {1 3}}, \phi_ {2} = \frac {3 \chi_ {1} - 2 \chi_ {2}}{\sqrt {3}}
$$
当测量 $\hat{A}$ 后得到 $a_1$ , 若再测量 $\hat{B}$ , 然后再测量 $\hat{A}$ , 问第二次测量 $\hat{A}$ 得到 $a_1$ 的概率是多少?

\vfill

\newpage
\vfill

\noindent
\textbf{2-42}\quad 设能量 $E$ 是三度简并的, 对应的 3 个波函数为 $\phi_1, \phi_2, \phi_3$ , 它们不归一, 相互之间也不正交. 试通过它们, 构造出 3 个相互正交、且归一的波函数.

\vfill

\newpage
\vfill

\noindent
\textbf{2-43}\quad 一体系由两个不同类型的中微子组成, 其哈密顿量 $\hat{H}$ 的本征态矢为 $|1\rangle$ 与 $|2\rangle$ , 相应的本征能量为 $E_1$ 与 $E_2$ , $E_1 < E_2$ . 已知电子中微子与 $\mu$ 子中微子的态矢分别为
$$
\left| e \right\rangle = \cos \theta | 1 \rangle + \sin \theta | 2 \rangle , \left| \mu \right\rangle = - \sin \theta | 1 \rangle + \cos \theta | 2 \rangle
$$
其中 $\theta$ 是混合角. 设体系在 $t = 0$ 时处于电子中微子态 $\left|e\right\rangle$ , 求

(1)任意 $t$ 时刻体系的态矢 $\left|\psi(t)\right\rangle$ ; 

(2)任意 $t$ 时刻体系处于基态 $\left|1\right\rangle$ 的概率; 

(3)任意 $t$ 时刻体系处于 $\mu$ 子中微子态 $\left|\mu\right\rangle$ 的概率; 

(4)何时体系又回到电子中微子态 $\left|e\right\rangle$ , 周期是什么?

\vfill

\newpage
\vfill

\noindent
\textbf{2-44}\quad 一维谐振子降算符 $a$ 与升算符 $a^{+}$ 的定义为
\begin{equation}\tag{1}
a = \sqrt {\frac {\mu \omega}{2 \hbar}} \left(x + \frac {\mathrm{i}}{\mu \omega} \hat {p}\right), a ^ {+} = \sqrt {\frac {\mu \omega}{2 \hbar}} \left(x - \frac {\mathrm{i}}{\mu \omega} \hat {p}\right)
\end{equation}
由 $a$ 与 $a^{+}$ 可以构成厄米算符 $\hat{N} = a^{+}a$ . 令 $\hat{N}$ 的本征值为 $n$ , 本征态为 $\left|n\right\rangle$ . 

(1) 计算对易关系: $[a, a^{+}], [a, \hat{N}], [a^{+}, \hat{N}]$ ; 

(2) 证明
\begin{equation}\tag{2}
a | n \rangle = \sqrt {n} | n - 1 \rangle
\end{equation}
\begin{equation}\tag{3}
a ^ {+} | n \rangle = \sqrt {n + 1} | n + 1 \rangle
\end{equation}

(3)将谐振子的哈密顿算符 $\hat{H} = \frac{\hat{p}^2}{2\mu} + \frac{1}{2}\mu \omega^2 x^2$ 用 $\hat{N} = a^+ a$ 表示；

(4)利用式(2)求出 $\hat{N}$ 的本征值 $n$ ，从而求出 $\hat{H}$ 的本征值 $E$ ；

(5)由式(3)不难得到
\begin{equation}\tag{4}
\left| n \right\rangle = \frac {1}{\sqrt {n !}} (a ^ {+}) ^ {n} \left| 0 \right\rangle
\end{equation}
写出它在 $x$ 表象中的表达式 $\psi_{n}(x) = \langle x|n\rangle$ . 利用公式(2), 给出 $\psi_{0}(x) = \langle x|0\rangle$ 满足的方程, 并求出 $\psi_{0}(x)$ .

\vfill

\newpage
\vfill

\noindent
\textbf{2-45}\quad 一维谐振子降算符 $a$ 的本征值为 $\alpha$ 的本征态 $|\alpha\rangle$ 称为谐振子的相干态. 设谐振子能量本征态为 $|n\rangle$ , 由公式
\begin{equation}\tag{1}
a \left| n \right\rangle = \sqrt {n} \left| n - 1 \right\rangle
\end{equation}
看出, 谐振子的基态 $|0\rangle$ 是 $a$ 的本征值 $\alpha = 0$ 的本征态. $a$ 的本征值 $\alpha \neq 0$ 的本征态可以表示成 $|n\rangle$ 的线性叠加:
\begin{equation}\tag{2}
\left| \alpha \right\rangle = \sum_ {n = 0} ^ {\infty} c _ {n} \left| n \right\rangle
\end{equation}
求出 $c_{n}$ , 从而得到谐振子的相干态 $\left|\alpha\right\rangle$ .

\vfill

\newpage
\vfill

\noindent
\textbf{2-46}\quad 设 $|n\rangle$ 为一维谐振子的能量本征态, 谐振子降算符a对 $|n\rangle$ 的作用为
\begin{equation}\tag{1}
a | n \rangle = \sqrt {n} | n - 1 \rangle
\end{equation}
证明

(1) 一维谐振子相干态
\begin{equation}\tag{2}
\left| \alpha \right\rangle = \mathrm{e} ^ {- \left| \alpha \right| ^ {2} / 2} \sum_ {n = 0} ^ {\infty} \frac {\alpha^ {n}}{\sqrt {n !}} \left| n \right\rangle
\end{equation}

是谐振子降算符 $a$ 的本征态, 本征值为 $\alpha$ . 

(2) $a$ 的不同本征值 $\beta$ 与 $\alpha$ 的本征态不正交
\begin{equation}\tag{3}
\langle \beta | \alpha \rangle \neq 0
\end{equation}

(3) 谐振子相干态(2)可以表示为
\begin{equation}\tag{4}
\left| \alpha \right\rangle = \mathrm{e} ^ {\alpha a ^ {+} - \alpha^ {*} a} \left| 0 \right\rangle
\end{equation}

其中 $a^{+}$ 为谐振子升算符.

\vfill

\newpage
\vfill

\noindent
\textbf{2-47}\quad 一维谐振子处于相干态 $|\alpha\rangle = e^{-|\alpha|^2 / 2}\sum_{n=0}^{\infty} \frac{\alpha^n}{\sqrt{n!}} |n\rangle$ ，其中 $|n\rangle$ 是谐振子哈密顿量 $\hat{H} = \frac{\hat{p}^2}{2\mu} + \frac{1}{2}\mu \omega^2 x^2$ 的本征态。谐振子相干态 $|\alpha\rangle$ 是谐振子算符 $a = \sqrt{\frac{\mu \omega}{2\hbar}} \left(x + \frac{i}{\mu \omega} \hat{p}\right)$ 的本征值为 $\alpha$ 的本征态，满足归一化条件 $\langle \alpha | \alpha \rangle = 1$ 。

(1) 求能量平均值 $\bar{H}$ ；

(2) 求 $\Delta x$ 与 $\Delta p$ ，表明在相干态上乘积 $\Delta x \Delta p$ 取测不准关系式 $\Delta x \Delta p > \hbar / 2$ 中的最小值。

\vfill

\newpage
\vfill

\noindent
\textbf{2-48}\quad 设厄米电荷算符 $\hat{Q}$ 的本征态为 $\left|\psi_q\right\rangle$ ，本征值为 $q: \hat{Q} \left|\psi_q\right\rangle = q \left|\psi_q\right\rangle$ 。电荷共轭算符 $\hat{C}$ 对 $\left|\psi_q\right\rangle$ 的作用是使之成为 $\hat{Q}$ 的本征值为 $-q$ 的本征态 $\left|\psi_{-q}\right\rangle$ ： $\hat{C} \left|\psi_q\right\rangle = \left|\psi_{-q}\right\rangle$ 。证明算符 $\hat{C}$ 与 $\hat{Q}$ 反对易： $\hat{C}\hat{Q} + \hat{Q}\hat{C} = 0$ 。

\vfill

\newpage